 \documentclass[10pt]{amsart}
\usepackage[OT1]{fontenc}
\usepackage[margin = 1in]{geometry}

%packages
\usepackage{amsfonts,amsmath,amssymb,amsthm}
\usepackage{enumerate}
\usepackage[mathscr]{euscript}
\usepackage{tikz-cd}
\usepackage{quiver}
\usepackage{wasysym}
\usepackage{stmaryrd}
\usepackage{subcaption}
\usepackage{bbm}

\usepackage{tabularx}
\usepackage{booktabs}
\usepackage[labelfont=sc,format=plain,justification=centering,singlelinecheck=false,tablename=Index of Notation.]{caption}


%bibliography & hyperlinks
\usepackage[backend=biber, useprefix=true, style = alphabetic]{biblatex}
\usepackage[dvipsnames]{xcolor}
\addbibresource{refs.bib}
\usepackage{hyperref}
\hypersetup{
    colorlinks = true,
    linktoc=page,
    urlcolor = BurntOrange,
    citecolor = Bittersweet,
    linkcolor = BurntOrange,
}

%\renewcommand{\emph}[1]{\textcolor{Regalia}{\textup{#1}}}



\numberwithin{equation}{section}

%define anything useful
\def\~#1{\mathscr{#1}}
\def\cal#1{\mathcal{#1}}
\def\To#1{\xrightarrow{#1}}

%commands{}
\newcommand{\vp}{\varphi}
\newcommand{\ve}{\varepsilon}
\newcommand{\del}{\partial}
\newcommand{\G}{\mathbb{G}}
\newcommand{\N}{\mathbb{N}}
\newcommand{\C}{\mathbb{C}}
\newcommand{\A}{\mathbb{A}}
\newcommand{\R}{\mathbb{R}}
\newcommand{\Q}{\mathbb{Q}}
\newcommand{\Z}{\mathbb{Z}}
\newcommand{\M}{\mathrm{M}}
\newcommand{\SL}{\mathrm{SL}}
\newcommand{\SO}{\mathrm{SO}}
\renewcommand{\O}{\mathrm{O}}
\newcommand{\isom}{\xrightarrow{\sim}}
\newcommand{\epi}{\twoheadrightarrow}
\newcommand{\mono}{\hookrightarrow}
\newcommand{\im}{\mathop{\mathrm{im}}}
\newcommand{\coker}{\mathop{\mathrm{coker}}}
\newcommand{\colim}{\mathop{\mathrm{colim}}}
\newcommand{\Char}{\mathrm{char}}
\newcommand{\lcm}{\mathrm{lcm}}
\newcommand{\ord}{\mathrm{ord}}
\newcommand{\Gcd}{\mathrm{gcd}}
\newcommand{\Aut}{\mathrm{Aut}}
\newcommand{\Stab}{\mathrm{Stab}}
\newcommand{\Spec}{\mathrm{Spec}}
\newcommand{\GL}{\mathrm{GL}}
\newcommand{\affGL}{\~G\~L}
\newcommand{\Grp}{\~G\mathrm{rp}}
\newcommand{\Gr}{\mathrm{Gr}}
\newcommand{\Coh}{\mathrm{Coh}}
\newcommand{\YMH}{\mathrm{YMH}}
\newcommand{\id}{\mathrm{id}}
\newcommand{\fg}{\mathrm{fg}}
\newcommand{\Proj}{\mathrm{Proj}}
\newcommand{\op}{\mathrm{op}}
\newcommand{\ev}{\mathrm{ev}}
\newcommand{\aff}{\mathrm{aff}}
\newcommand{\fd}{\mathrm{fd}}
\newcommand{\ob}{\mathrm{ob}}
\newcommand{\rk}{\mathrm{rk}}
\newcommand{\rank}{\mathrm{rank}}
\renewcommand{\div}{\mathrm{div}}
\newcommand{\Hom}{\mathrm{Hom}}
\newcommand{\iHom}{\underline{\Hom}}
\newcommand{\Div}{\mathrm{Div}}
\newcommand{\PDiv}{\mathrm{PDiv}}
\DeclareMathOperator{\End}{End}
\DeclareMathOperator{\tr}{tr}
\newcommand{\QCoh}{\mathrm{QCoh}}
\newcommand{\QC}{\mathrm{QC}}

\newcommand{\Pic}{\mathrm{Pic}}
\renewcommand{\frak}[1]{\mathfrak{#1}}
\newcommand{\Var}{\mathrm{Var}}
\renewcommand{\Vec}{\mathrm{Vec}}
\newcommand{\Vect}{\mathrm{Vect}}
\newcommand{\Rep}{\mathrm{Rep}}
\newcommand{\Fun}{\mathrm{Fun}}
\newcommand{\Hopf}{\mathrm{Hopf}}
\newcommand{\gSch}{\mathrm{gSch}}
\newcommand{\Mod}{\mathrm{Mod}}
\newcommand{\Alg}{\mathrm{Alg}}
\newcommand{\Aff}{\mathrm{Aff}}

\renewcommand{\sectionautorefname}{\S \!}
\renewcommand{\subsectionautorefname}{\S \!}
\renewcommand{\subsubsectionautorefname}{\S\!}

%theorem environments
\theoremstyle{plain}
\newtheorem{thm}{Theorem}[section]
\newtheorem*{thm*}{Theorem}
\newtheorem{prop}[thm]{Proposition}
\newtheorem{lem}[thm]{Lemma}
\newtheorem*{lem*}{Lemma}
\newtheorem{cor}[thm]{Corollary}
\newtheorem*{cor*}{Corollary}

\theoremstyle{definition}
\newtheorem{defn}[thm]{Definition}
\newtheorem*{defn*}{Definition}
\newtheorem{ex}[thm]{Example}
\AtEndEnvironment{ex}{\null\hfill$\Diamond$}
\newtheorem{exs}[thm]{Examples}
\AtEndEnvironment{exs}{\null\hfill$\Diamond$}

\theoremstyle{remark}
\newtheorem*{rmk}{Remark}
\newtheorem*{warn}{Warning}
\newtheorem*{rmks}{Remarks}
\newtheorem*{notation}{Notation}

\begin{document}

\title{Tannaka duality and the fundamental group scheme}


\author{Elias Vandenberg}
\address{Department of Mathematics, Western University}
\email{evande95@uwo.ca}
\date{\today}

\maketitle

\tableofcontents
%\section*{Introduction}


\subsection*{Conventions and Notation}
All rings and algebras will be assumed to be commutative with identity $1 \ne 0$. The symbol $:=$ indicates that the definition in question is an equality. If $\~C$ is a category and $X$ is an object of $\~C$, we write $X \in \~C$ in place of $X \in \ob(\~C)$. We write $\~C(X,Y)$ the collection of morphisms $X \to Y$ in $\~C$ (when working with categories of groups/modules/etc, we instead use the notation $\Hom(X,Y)$). $\Fun(\~C, \~D)$ denotes the category of functors $\~C \to \~D$. The symbol $\cong$ denotes an isomorphism, while $\simeq$ is reserved for equivalences of categories.



\section{Tannaka's theorem for linear algebraic groups} In this section, $k$ denotes a field. Tensor products without subscripts are over $k$.

\subsection{Review of linear algebraic groups}

\begin{defn}
    Let $\~C$ be a category with a terminal object $I$. A \emph{group internal to $\~C$} (or simply a \emph{group in} $\~C$) is a tuple $(G, c, i, e)$ consists of the following data:
    \begin{enumerate}
        \item An object $G \in \~C$ such that finite products of $G$ with itself exist.
        \item Maps $c: G \times G \to G$ (\emph{composition}), $i: G \to G$ (\emph{inversion}) and $e: I \to G$ (\emph{identity}) such that the following diagrams commute:
\begin{equation}\label{eq:1.1}
\begin{tikzcd}[sep=scriptsize]
    {G \times G \times G} & {G \times G } && {G \times G} & {G \times G} && G & {G\times G} \\
    &&& G & G \\
    {G \times G} & G && {G \times G} & {G \times G} && {G \times G} & G
    \arrow["{{c \times \id}}", from=1-1, to=1-2]
    \arrow["{{\id \times c}}"', from=1-1, to=3-1]
    \arrow["c", from=1-2, to=3-2]
    \arrow["c"', from=1-4, to=2-4]
    \arrow["{{\id \times i}}"', from=1-5, to=1-4]
    \arrow["{{\id \times e}}", from=1-7, to=1-8]
    \arrow["{{e \times \id}}"', from=1-7, to=3-7]
    \arrow[equals, from=1-7, to=3-8]
    \arrow["c", from=1-8, to=3-8]
    \arrow["{{\langle \id, \id \rangle}}"', from=2-5, to=1-5]
    \arrow["{{1_G \mapsfrom g }}"', from=2-5, to=2-4]
    \arrow["{{\langle \id, \id \rangle}}", from=2-5, to=3-5]
    \arrow["c"', from=3-1, to=3-2]
    \arrow["c", from=3-4, to=2-4]
    \arrow["{{i \times \id}}"', from=3-5, to=3-4]
    \arrow["c"', from=3-7, to=3-8]
\end{tikzcd}\end{equation}
(here we identify $G\times(G \times G)$ with $(G \times G) \times G$ in the obvious manner). The three diagrams are called the associativity, inverse, and identity axioms respectively.
    \end{enumerate}
\end{defn}

\begin{defn}
    An \emph{affine group scheme} is an affine scheme $G = \Spec(A)$ which is a group object in $\Aff_k$. When $G$ is an algebraic variety (i.e., $A$ is finitely-generated), we call $G$ an \emph{algebraic group}, and if $G$ is an affine variety, we call $G$ an \emph{affine algebraic group} or a \emph{linear algebraic group}. We denote by $\gSch_k^\aff$ the category of affine group schemes over $k$ 
\end{defn}
Consider a linear algebraic group $G = \Spec(A)$. To understand the group structure on $G$, we need to understand the group structure on the underlying $k$-algebra $A$. $\Spec$ reverses arrows, and the product of two affine varieties corresponds to the tensor product of their associated coordinate rings, so the morphisms  $c: G \times G \to G$, $i: G \to G$, and $e: I \to G$ induce maps, $\Delta: A \to A \otimes_k A$ (called \emph{comultiplication}), $\iota: A \to A$ (called the \emph{coinverse} or \emph{antipode}), and $\ve: A \to k$ (the \emph{counit}). Furthermore, the composite map $\eta: A \To{\ve} k \mono A$ corresponds to the morphism $g \mapsto 1_G$, and we have a multiplication map $\mu: A \otimes_k A \to k$ given by $\mu(a \otimes b) = ab$. Reversing the arrows in (\ref{eq:1.1}), we see that the following diagrams must commute:
\begin{equation}\label{eq:1.2}\begin{tikzcd}[sep = scriptsize]
    A & {A \otimes A} && {A \otimes A} & {A \otimes A} && A & {A \otimes A} \\
    &&& A & A \\
    {A \otimes A} & {A \otimes A \otimes A} && {A \otimes A} & {A \otimes A} && {A \otimes A} & A
    \arrow["\Delta", from=1-1, to=1-2]
    \arrow["\Delta"', from=1-1, to=3-1]
    \arrow["{\id \otimes \Delta}", from=1-2, to=3-2]
    \arrow["{\iota \otimes \id}", from=1-4, to=1-5]
    \arrow["\mu", from=1-5, to=2-5]
    \arrow["{\ve \otimes \id}"', from=1-8, to=1-7]
    \arrow["\Delta", from=2-4, to=1-4]
    \arrow["\eta", from=2-4, to=2-5]
    \arrow["\Delta"', from=2-4, to=3-4]
    \arrow["{\Delta \otimes \id}"', from=3-1, to=3-2]
    \arrow["{\id \otimes \iota}"', from=3-4, to=3-5]
    \arrow["\mu"', from=3-5, to=2-5]
    \arrow["{\id \otimes \ve}", from=3-7, to=1-7]
    \arrow[equals, from=3-8, to=1-7]
    \arrow["\Delta"', from=3-8, to=1-8]
    \arrow["\Delta", from=3-8, to=3-7]
\end{tikzcd}\end{equation}
These are called the coassociativity, coinverse, and counit axioms respectively. This gives $(A, \Delta, \iota, \ve)$ the structure of a commutative Hopf algebra over $k$. Hence every affine group scheme $G$ over $k$ corresponds to $\Spec(A)$ for some commutative Hopf algebra $A$ over $k$, so there is an equivalence of categories $\gSch_k^\aff \simeq \left(\Hopf_k\right)^\op$ (where $\Hopf_k$ denotes the category of commutative Hopf algebras over $k$).
\begin{ex}We record a few important examples of linear algebraic groups.
 \begin{enumerate}
\item Take $G = \A^1 = k$ with the group operation given by addition, and let $A = k[t]$ be the associated coordinate ring. The algebraic variety $G$ then has a group structure with 
\[c(x, y) = x + y, \quad i(x) = -x, \quad \text{and} \quad e(x) = 0.\]
The comultiplication $\Delta: A \to A \otimes_k A$ is given $\Delta(t) = t \otimes 1 + 1 \otimes t$, the coinverse $\iota: A \to A$ is given by $\iota(t) = -t$, and the counit $\ve: A \to k$ is given by $\ve(t) = 0$. We refer to $G$ as the \emph{additive group}, denoted by $\G_a$.
\item Similarly, we may take $G = k^\times$, where now the group operation is given by multiplication. In this case, we have $A = k[t, t^{-1}]$. The comultiplication $\Delta: A \to A \otimes_k A$ is given by $\Delta(t) = t \otimes t$, the coinverse $\iota: A \to A$ is given by $\iota(t) = t^{-1}$, and the counit $\ve: A \to k$ is given by $\ve(t) = 1$. We refer to $G$ as the \emph{multiplicative group}, denoted by $\G_m$.
\item         Let $\M_n(k)$ denote the set of $n \times n$ matrices over $k$, which we identify with $k^{n^2}$ in the obvious manner. Recall that the \emph{general linear group} $\GL_n(k)$ is then defined as
\[\GL_n(k) := \{X \in \M_n(k) : \det(X) \ne 0\},\]
where the group operation is given by matrix multiplication. $\GL_n(k)$ is an algebraic variety with coordinate ring 
\[A = k[\GL_n] = k[X_{ij}, \det(X)^{-1}]_{1 \le i, j \le n}\]
where $X = (X_{ij})_{1 \le i, j \le n}$. The comultiplication $\Delta: A \to A \otimes_k A$ is given by 
\[\Delta(X_{ij}) = \sum_{k =1}^n X_{ik} \otimes X_{kj}, \quad 1 \le i, j \le n,\]
$\iota: A \to A$ is given by $X_{ij} \mapsto X_{ij}^{-1}$, the $(i,j)$ entry of the matrix $X^{-1}$, and $\ve(X_{ij}) = \delta_{ij}$ (the Kronecker delta). More generally, if $V$ is a $k$-vector space of dimension $n < \infty$, $\GL(V)$ has the structure of an algebraic group.   
\item     Any subgroup of $\GL_n(k)$ which is closed in the Zariski topology is a linear algebraic group. Well-known examples include:
    \begin{enumerate}
        \item Any finite subgroup of $\GL_n(k)$.
        \item The special linear group $\SL_n(k) = \{X \in \GL_n(k) : \det(X) = 1\}$.
        \item The orthogonal group $\O_n(k) = \{X \in \GL_n(k) : X^T X = I\}$.
        \item The special orthogonal group $\SO_n(k) = \O_n(k) \cap \SL_n(k)$.
    \end{enumerate} 
\end{enumerate}
\end{ex}

\subsection{Representations of linear algebraic groups}

Throughout this section, we fix an affine group scheme $G$ over $k$.
\begin{defn}
A \emph{representation} of $G$ is pair $(V, r_V)$ where $V$ is a vector space $V$ (called a \emph{$G$-module}) and $r_V: G \to \GL(V)$ is a morphism of affine group schemes. We write $I = k$ for the trivial $G$-module ($r_I(g) = 1$ for all $g \in G$).    Given $G$-modules $(V, r_V)$ and $(W, r_W)$, a \emph{morphism} or \emph{$G$-equivariant map} $(V, r_V) \to (W, r_W)$ is a morphism $\vp: V \to W$ such that the following diagram commutes for every $g \in G$:
    \begin{equation}\label{eq:1.3}
    \begin{tikzcd}[sep=scriptsize]
    V & W \\
    V & W
    \arrow["\vp", from=1-1, to=1-2]
    \arrow["{r_V(g)}"', from=1-1, to=2-1]
    \arrow["{r_W(g)}", from=1-2, to=2-2]
    \arrow["\vp"', from=2-1, to=2-2]
\end{tikzcd}\end{equation}
We say that a representation $(r_V, V)$ is \emph{finite-dimensional} when $V$ is finite-dimensional.

\end{defn}

We can rephrase the above definitions in terms of categories: Regard the affine group scheme $G$ as a one object groupoid (so each element $g \in G$ corresponds to an invertible morphism $g: G \to G$), and consider the category $\Vec_k$ of finite-dimensional vector spaces over a field $k$. A functor $r_V: G \to \Vec_k$ consists of a single vector space $V$ (the value of $r_V$ at the unique object $G$), together with a $k$-linear map $r_V(g): V \to V$ for each $g \in G$. Therefore a functor $r_V: G \to \Vec_k$ is an assignment of each $g \in G$ to a $k$-linear map in $\End(V) = \GL(V)$ which respects composition (the group operation) and identity. In other words, such a functor is simply a group homomorphism $G \to \GL(V)$, i.e., a linear representation of $G$. If $(V, r_V)$ and $(W, r_W)$ are linear representations of $G$, a morphism $(V, r_V) \to (W, r_W)$ is simply a natural transformation between the associated functors $r_V$ and $r_W$. This consists of a single $k$-linear map $\vp:  V \to W $, such that for every $g \in G$ the diagram (\ref{eq:1.3}) commutes. Thus, an element of $\Fun(G, \Vec_k)$ is a representation of $G$, and a morphism in this category is precisely a $G$-equivariant map. Given an affine group scheme $G$ and a $k$-vector space $V$, we define a category of all $k$-linear representations and equivariant maps by
\[\Rep'(G) := \Fun\left(G, \Vec_k \right).\] 
We therefore have a forgetful functor $\omega_G': \Rep'(G) \to \Vec_k$, defined on objects by $(V, r_V) \mapsto V$ and on morphisms by $(V \To{f} W) \mapsto (V \To{f} W)$ (now viewed only as a morphism of vector spaces). We will almost always work exclusively with finite-dimensional representations of group schemes. We write
\[\Rep(G) := \Fun\left(G, \Vec_k^\fd\right)\]
for the category of finite-dimensional representations of $G$. This again comes with a forgetful functor $\omega_G: \Rep(G) \to \Vec_k^\fd$. 

If $V$ is a finite-dimensional $k$-vector space, recall that the \emph{dual space} of $V$ is the vector space 
\[V^\vee := \Hom(V, k).\]
There is a natural \emph{duality pairing} $\langle \cdot, \cdot \rangle: V \times V^\vee \to V$ given by
\[\langle v, u \rangle := u(v) \quad \text{for all } v \in V, \, u \in V^\vee.\]
If $V$ is a $G$-module, then $V^\vee$ is also a $G$-module; the map $r_{V^\vee}: G \to \GL(V^\vee)$ is given by
\[\left(r_{V^\vee}(g)(u)\right)v := u \left(r_V(g^{-1})(v) \right) \quad \text{for all } u \in V^\vee.\]
This pairing satisfies the following invariance property:
\[\langle r_V(g)(v), r_{V^\vee}(g)(u) \rangle = \langle v, u \rangle \quad \text{for all } v \in V, \, u \in V^\vee.\]

When studying $G$-modules and their representations, it is natural to ask how $G$ acts on its coordinate ring $k[G]$. However, this presents some problems: $k[G]$ is certainly finite-dimensional as a $k$-algebra, but it may not be finite-dimensional as a $k$-vector space. To remedy this issue, we will have to enlarge the category of $G$-modules that we are working in. This leads to the notion of a \emph{locally finite $G$-module}.

\begin{defn}
    Let $V$ be any vector space (not necessarily finite-dimensional). We say that $\alpha \in \End(V)$ is \emph{locally finite} if $V$ is a union of finite-dimensional subspaces, each of which is stable under $\alpha$. We say that the $G$-module $\rho: G \to \GL(V) = \End(V)$ is locally finite if $\rho(g)$ is locally finite for all $g \in G$. Hence $V$ is locally finite if it can be written as a union of finite-dimensional $G$-stable subspaces. 
\end{defn}
\begin{ex}
\hfill
\begin{enumerate}
\item The \emph{trivial representation} $(\mathbf{1}, r_\mathbf{1})$ is defined by taking $\mathbf{1} = k$ and $r_\mathbf{1}(g) = \id_k$ for all $g \in G$.
\item     Given representations $(V, r_V)$ and $(W, r_W)$, the tensor product $V \otimes W$ has the structure of a $G$-module, with $r_{V \otimes W} = r_V \otimes r_W$.   
\end{enumerate} 
\end{ex}
\begin{ex}\label{ex:1.5}
    Let $G$ be an algebraic group. Then $k[G]$ is a $G$-module: the action by $g \in G$ is given by
    \begin{align*}
        \rho(g): k[G] &\to k[G]\\
        F &\mapsto (g\cdot F)(h)
    \end{align*}
    where 
    \[(g \cdot F)(h) = F(hg)\]
    for all $F \in k[G]$. 
\end{ex}

\begin{prop}
    Given an algebraic group $G$, the action above endows $A$ with the structure of a locally finite $G$-module. 
\end{prop}


\begin{defn}
    A \emph{coalgebra} over $k$ is a $k$-vector space $C$ equipped with $k$-linear maps $\Delta: C \to C \otimes_k C$ and $\ve: C \to k$ satisfying the coassociativity and counit axioms in (\ref{eq:1.2}). A \emph{comodule} over a coalgebra $C$ is a $k$-vector space $V$ with a $k$-linear map $\rho: V \to V \otimes_k C$ such that the following diagrams commute:
    \begin{equation}\begin{tikzcd}
    V & {V \otimes C} && V & {V \otimes C} \\
    & {V\otimes k \cong V} && {V \otimes C} & {V \otimes C \otimes C}
    \arrow["\rho", from=1-1, to=1-2]
    \arrow[equals, from=1-1, to=2-2]
    \arrow["{\id \otimes \ve}", from=1-2, to=2-2]
    \arrow["\rho", from=1-4, to=1-5]
    \arrow["\rho"', from=1-4, to=2-4]
    \arrow["{\rho \otimes \id}", from=1-5, to=2-5]
    \arrow["{\id \otimes \Delta }"', from=2-4, to=2-5]
\end{tikzcd}\end{equation}
Notice that $\Delta$ defines a $C$-comodule structure on $C$.
\end{defn}

\begin{prop}[\cite{deligne1982tannakian}, Proposition 2.2]\label{prop:DM2.2}
    For any affine group scheme $G$ over $k$ and every $k$-vector space $V$, there is a one-to-one correspondence between $A$-comodule structures on $V$ and linear representations of $G$ on $V$.
\end{prop}
\begin{proof}

\end{proof}
\begin{defn}
    The representation of $G$ on $A$ defined by $A$'s comodule structure $\Delta$ is called the \emph{regular representation}.
\end{defn}

\subsection{Proof of Tannaka's theorem} 
Given $G$-modules $V$ and $W$, recall that $V \otimes W$ is a $G$-module with $r_{V \otimes W} = r_V \otimes r_W$. We will need the following lemma:
\begin{lem}\label{lem:1}
    Let $G = \Spec(A)$ be a linear algebraic group and $V$ a $G$-module. For $v \in V$, $u \in V^\vee$, and $g \in G$, define $\vp_u(v)(g) = \langle r_V(g)v, u \rangle$. Then $\vp_u(v) \in A$ and $\vp_u: V \to A$ is a morphism of $G$-modules.
\end{lem}

\begin{proof}
    To prove the first assertion, one simply notes that $\vp_u(v)$ is given by the composite map
    \[\vp_u(v): \begin{tikzcd}
    G & {\GL(V)} && V & {k = \A^1}
    \arrow["{r_V}", from=1-1, to=1-2]
    \arrow["{v \mapsto r_V(g)(v)}", from=1-2, to=1-4]
    \arrow["u", from=1-4, to=1-5]
\end{tikzcd}\]
The maps $r_V$ and $v \mapsto r_V(g)(v)$ are morphisms, and $u$ is a morphism, since after choosing a basis $\{e_1, \hdots, e_n\}$ for $V$, we can write $u(e_1, \hdots, e_n) = c_1e_1 + \cdots + c_n e_n$. Hence the composite $\vp_u(v)$ defines a morphism of varieties $G \to \A^1$, so $\vp_u(v) \in A$. 

We now check that $\vp_u: V \to A$ is a morphism of $G$-modules, i.e., the following diagram commutes for all $g \in G$:
\[\begin{tikzcd}[sep=scriptsize]
    V & A \\
    V & A
    \arrow["\vp_u", from=1-1, to=1-2]
    \arrow["{r_V(g)}"', from=1-1, to=2-1]
    \arrow["{\rho(g)}", from=1-2, to=2-2]
    \arrow["\vp_u"', from=2-1, to=2-2]
\end{tikzcd}\]
We must show that $\rho(g)(\vp_u(v)) = \vp_(u_V(g))$ in $A$. The elements of $A = k[G]$ are functions on $G$, so we need to show that
\[\left(\rho(g)(\vp_u(v)) \right)(h) = \left(\vp_u(r_V(g))\right)(h) \quad \text{for all } h \in G.\]
Given $h \in G$, the left-hand side becomes
\[\left(\rho(g)(\vp_u(v)) \right)(h) = \vp_u(v)(hg),\]
and on the right-hand side, we have
\begin{align*}
    \left(\vp_u(r_V(g))\right)(h) &= \langle r_V(h)(r_V(g)v), u \rangle \\ 
                                  &= \langle r_V(hg)v, u \rangle \\ 
                                  &= \vp_u(v)(hg),
\end{align*}
so we see that the diagram commutes.
\end{proof}

\begin{thm}[Tannaka]\label{thm:1}
    Let $V$ be a finite-dimensional $G$-module, and suppose there exists an element $\alpha_V \in \GL(V)$ with the following properties:
    \begin{enumerate}
        \item \label{T1} If $V$ and $W$ are $G$-modules, then $\alpha_{V \otimes W} = \alpha_V \otimes \alpha_W$.
        \item \label{T2} If $\psi: V \to W$ is a homomorphism of $G$-modules, then $\psi \circ \alpha_V = \alpha_W \circ \psi$.
        \item \label{T3} $\alpha_{\mathbf{1}} = 1$.
    \end{enumerate} 
    Then there exists $x \in G$ such that $\alpha_V = r_V(x)$ for all $G$-modules $V$.
\end{thm}

\begin{proof}
We begin by defining $\alpha_V$ for a locally finite $G$-module $V$. We will then use this construction to prove the theorem for a finite-dimensional $G$-module. Let $v \in V$ and let $W \subseteq V$ be a finite dimensional subspace of $V$ which contains $v$. Define $\alpha_V(v) = \alpha_W(v)$. If $W_2 \subseteq V$ is another finite-dimensional $G$-stable subspace containing $V$, set $W_{12} = W_1 \cap W_2$, and denote by $i_1: W_{12} \mono W_1$ and $i_2: W_{12} \mono W_2$ the inclusions. Take $w \in W_1 \cap W_2$. Applying (2) to these inclusions, we get 
\[\left(i_1 \circ \alpha_{W_{12}}\right)(w) = \alpha_{W_{12}}(w) = \alpha_{W_1}(w) = (\alpha_{W_1} \circ i_1)(w),\]
and
\[\left(i_2 \circ \alpha_{W_{12}}\right)(w) = \alpha_{W_{12}}(w) = \alpha_{W_2}(w) = (\alpha_{W_2} \circ i_2)(w),\]
so $\alpha_{W_1} = \alpha_{W_{12}} = \alpha_{W_2}$ and we see that the map is well-defined. It is clearly an invertible linear map, since each $\alpha_W$ is invertible and linear. Now suppose $V$ and $V'$ are both locally finite. Then we can find a dimensional subspace $W \subseteq V$ with $v \in W$. As $W$ and $W'$ are both finite-dimensional $G$-modules, (1) gives $\alpha_{W \otimes W'} = \alpha_W \otimes \alpha_{W'}$. Since $\alpha$ is defined on finite-dimensional subspaces, we therefore have $\alpha_{V \otimes V'} = \alpha_V \otimes \alpha_{V'}$.  Moreover, if $\psi: V \to V'$ is a homomorphism of $G$-modules, then $W' = \psi(V)$ is a finite-dimensional $G$-stable subspace of $V$, so the restriction $\psi|
_{W}: W \to W'$ is $G$-equivariant, hence $\psi|_{W} \circ \alpha_W = \alpha_{W'} \circ \psi|_W$ by (2). This gives 
\begin{align*}
    \left(\psi \circ \alpha_V\right)(v) &=  \left(\psi|_{W} \circ \alpha_W \right)(v)\\ 
                                     &= (\alpha_{W'}\circ \psi|_W)(v)\\
                                     &= \left(\alpha_{V'} \circ \psi\right)(v)
\end{align*}
so the map $\alpha_V$ satisfies properties (1) and (2).

Now consider the locally finite $G$-module $A = k[G]$ with the representation $\rho$ of Example \ref{ex:1.5} (with the action on $A \otimes A$ being given by $\rho \otimes \rho$). The diagram
\[\begin{tikzcd}
    {A \otimes A} & A \\
    {A \otimes A} & A
    \arrow["\mu", from=1-1, to=1-2]
    \arrow["{(\rho \otimes\rho)(g)}"', from=1-1, to=2-1]
    \arrow["{\rho(g)}", from=1-2, to=2-2]
    \arrow["\mu"', from=2-1, to=2-2]
\end{tikzcd}\]
is easily seen to be commutative, so the multiplication map $\mu: A \otimes A \to A$ is $G$-equivariant. Now consider $\alpha_A: A \to A$. We have $\alpha_{A \otimes A} = \alpha_A \otimes \alpha_A$ by (1), and applying (2) to the map $\mu$ gives $\mu \circ (\alpha_A \otimes \alpha_A) = \alpha_A \circ \mu$. We showed above that $\alpha_V$ is a well-defined automorphism of $V$ for any locally-finite $G$-module, and $\alpha_A$ commutes with $\mu$, so $\alpha_A$ is an automorphism of the $k$-algebra $A$. Applying $\Spec$ gives an automorphism $\theta: G \to G$ of $G = \Spec(A)$ such that $(\alpha F)(g) = F(\theta(g))$ for all $F \in A$.

Now consider the homomorphism $\Delta: A \to A \otimes A$. The image of $F \in A$ under $\Delta$ is a new map $\Delta F \in A \otimes A$, defined on $G \times G$ by $\Delta F(g, h) = F(gh)$. We claim that $\Delta \circ \rho(g) = \left(\id \otimes \rho(g)\right) \circ \Delta$, i.e., the following diagram commutes:
\[\begin{tikzcd}
    A & {A \otimes A} \\
    A & {A \otimes A}
    \arrow["\Delta", from=1-1, to=1-2]
    \arrow["{\rho(g)}"', from=1-1, to=2-1]
    \arrow["{\id \otimes \rho(g)}", from=1-2, to=2-2]
    \arrow["\Delta"', from=2-1, to=2-2]
\end{tikzcd}\]
Indeed, for all $F' \in A \otimes A$ and all $(h, t) \in G \times G$, we have $((\id \otimes \rho(g))F')(h,t) = F'(h, tg)$. Applying this to $\Delta F$ gives
\begin{align*}
   ((\id \otimes \rho(g))\Delta F)(h,t) &= \Delta F(h, tg) \\ 
                                        &= f(htg) \\ 
                                        &= \rho(g) F(ht) \\ 
                                        &= \Delta(\rho(g)F)(h,t),
\end{align*}
so the diagram commutes. Let $B = A \otimes A$ with the representation $r_B = \id \otimes \rho$. Then $\Delta: A \to B$ is equivariant by our previous calculation, so (1), (2), and (3) give $\alpha_B = \alpha_{A \otimes A} = \id \otimes \rho$. As $\Delta$ is a homomorphism, (2) gives $\Delta \circ \alpha_A = \left(\id \otimes \alpha_A\right) \circ \Delta$. Recalling $(\alpha_A F)(g) = F(\theta g)$ for all $F \in A$, this condition becomes $\theta(gh) = g \theta(h)$ for all $F \in A$, $g \in G$. We set $x = \theta(e)$, so $\theta(g) = gx$ and $\alpha_A = \rho(x)$ for all $x \in G$.

Now let $V$ be a finite-dimensional $G$-module. Given $u \in V^\vee$, we have the $G$-morphism $\vp_u: V \to A$ of Lemma \ref{lem:1}. Applying property (2) gives 
\[\vp_u \circ \alpha_V = \alpha_A \circ \vp_u.\]
Noting that $(\alpha_A \circ \vp_u)(g) = \vp_u(\theta(g)) = \vp_u(gx)$, the equation above becomes
\[\langle r_V(g) \alpha_V(v), u \rangle = \langle r_V(gx)v, u\rangle,\]
and this is true for any $u \in V^\vee$, so for all $g \in G$ we have
\[r_V(g) \alpha_V(v) = r_V(gx)v.\]
In particular, taking $g = 1$, we get $\alpha_V(v) = r_V(x)v$ for all $v \in V$, so $\alpha_V = r_V(x)$, as desired. 
\end{proof}

Let $G$ be a linear algebraic group over the field $k$. Recall that we have a forgetful functor
\[\omega_G: \Rep(G) \to \Vec_k^{\fd}.\] 
Tannaka's theorem is a statement about the natural isomorphisms of this functor. An element here is a collection of automorphisms 
\[\alpha = \{\alpha_V: \omega_G(V) \to \omega_G(V)\}_{(V, r_V) \in \Rep(G)}\]
such that for every $G$-equivariant map $f: V \to W$, the following diagram commutes:
\[\begin{tikzcd}
    {\omega_G(V)} & {\omega_G(W)} \\
    {\omega_G(V)} & {\omega_G(W)}
    \arrow["f", from=1-1, to=1-2]
    \arrow["{\alpha_V}"', from=1-1, to=2-1]
    \arrow["{\alpha_W}", from=1-2, to=2-2]
    \arrow["f"', from=2-1, to=2-2]
\end{tikzcd}\]

A ``good'' automorphism of $\omega_G$ should respect this structure in the sense of the conditions in Tannaka's theorem. Call such an automorphism a \emph{tensor automorphism}, and denote the collection of tensor automorphisms by $\underline{\Aut}^\otimes (\omega_G)$. For each $g \in G$, define a natural transformation $\alpha^g$ as follows:
\[\alpha_V^g: \omega_G(V) \to \omega_G(V), \quad v \mapsto r_V(g)(v). \]
This defines a group homomorphism
\[G \to \underline{\Aut}^\otimes(\omega_G), \quad g \mapsto \alpha^g.\]
Tannaka's theorem asserts that this map is actually an isomorphism. Hence every automorphism of the forgetful functor $F_G$ satisfying the three properties in Theorem \ref{thm:1} corresponds to the action of some element in $G$.
\section{Tannakian categories}

\subsection{Tensor categories}\label{sub:2.1}
In this section, we generalize Tannaka's theorem to the case of an affine group scheme. We adopt the terminology of \cite{deligne1982tannakian}, meaning that we will speak of \emph{tensor categories} as opposed to symmetric monoidal categories. This distinction in language will become more relevant when we introduce Tannakian categories, which contain much more algebraic data than the average symmetric monoidal category.

\subsubsection{Definition of tensor categories}\label{ssub:2.1.1}
Let $\~C$ be a category equipped with a functor $\otimes: \~C \times \~C \to \~C$ (written $\otimes(X, Y) = X \otimes Y$). An \emph{associativity constraint} for $\~C$ is a natural isomorphism $\alpha$ with components given by
\[\alpha_{X, Y, Z}: X \otimes (Y \otimes Z) \to (X \otimes Y) \otimes Z,\]
such that for all $W, X, Y, Z \in \~C$, the following diagram commutes:
\begin{equation}\label{eq:pentagon}
\begin{tikzcd}[column sep=tiny,row sep=scriptsize]
    & {W \otimes(X \otimes(Y \otimes Z))} \\
    {W \otimes ((X \otimes Y) \otimes Z)} && {(W \otimes X) \otimes (Y \otimes Z)} \\
    \\
    {(W \otimes (X \otimes Y)) \otimes Z} && {((W \otimes X) \otimes Y)\otimes Z}
    \arrow["{1 \otimes \alpha}"', from=1-2, to=2-1]
    \arrow["\alpha", from=1-2, to=2-3]
    \arrow["\alpha"', from=2-1, to=4-1]
    \arrow["\alpha", from=2-3, to=4-3]
    \arrow["{\alpha \otimes 1}"', from=4-1, to=4-3]
\end{tikzcd}
\end{equation}
This is known as the \emph{pentagon axiom}. A \emph{commutativity constraint} for $\~C$ is a natural isomorphism 
\[\beta_{XY}: X \otimes Y \to X \otimes Y\]
such that for all objects $X$ and $Y$, $\beta_{YX} \circ \beta_{XY}: X \otimes Y \to X \otimes Y$ is the identity on $X \otimes Y$:
\[\beta_{YX} \circ \beta_{XY} = 1_{X \otimes Y}.\]
Given associativity and commutativity constraints $\alpha$ and $\beta$ respectively, we say that $\alpha$ and $\beta$ are \emph{compatible} if the following diagram commutes for all $X, Y, Z \in \~C$:
\[\begin{tikzcd}[column sep=scriptsize]
    & {X \otimes(Y \otimes Z)} & {(X \otimes Y) \otimes Z} \\
    {X \otimes(Z \otimes Y)} &&& {Z \otimes (X \otimes Y)} \\
    & {(X \otimes Z) \otimes Y} & {(Z \otimes X) \otimes Y}
    \arrow["\alpha", from=1-2, to=1-3]
    \arrow["{1 \otimes \beta}"', from=1-2, to=2-1]
    \arrow["\beta", from=1-3, to=2-4]
    \arrow["\alpha"', from=2-1, to=3-2]
    \arrow["\alpha", from=2-4, to=3-3]
    \arrow["{\beta \otimes 1}"', from=3-2, to=3-3]
\end{tikzcd}\]
This condition is referred to as the \emph{hexagon axiom}. If $U \in \~C$ and $u: U \to U \otimes U$ is an isomorphism, the pair $(U, u)$ is an \emph{identity object} of $(\~C, \otimes)$ if the functor $\~C \to \~C$ given by $X \mapsto X \otimes U$ is an equivalence of categories.

\begin{defn}
    A \emph{tensor category}\footnote{In the language of category theory, a tensor category is a symmetric monoidal category.} is a tuple $(\~C, \otimes, \alpha, \beta)$ (or simply $(\~C, \otimes)$) where $\~C$ is a category with an identity object and $\alpha$ and $\beta$ are compatible associativity and commutativity constraints.
\end{defn}

\begin{ex}
    The category $\Mod_R$ of finitely-generated modules over a commutative ring $R$ is a tensor category with respect to the usual tensor product $M \otimes_R N$ (satisfying the obvious associativity and commutativity constraints). An identity object in this category is a pair $(U, u)$ where $U = Ru_0$ is a free $R$-module of rank $1$ with generator $u_0$, and $u: U \to U \otimes_R U$ is the unique isomorphism given by $u_0 \mapsto u_0 \otimes u_0$. Any generator $u_0$ determines such an identity object, and any identity object arises in this way. 
\end{ex}

\begin{prop}\label{prop:LeftUnitor}
    Let $(U, u)$ be an identity object in a tensor category $(\~C, \otimes)$.
    \begin{enumerate}
        \item There is a unique natural isomorphism (called the \emph{left unitor}), defined for $X \in \~C$ by
        \[\lambda_X: X \to U \otimes X,\]
        such that $\lambda_U = U$ and the following diagrams commute:
        \[\begin{tikzcd}[]
            {X \otimes Y} & {U \otimes (X \otimes Y)} && {X \otimes Y} & {(U \otimes X) \otimes Y} \\
            & {(U \otimes Y)\otimes Y} && {X \otimes (U \otimes Y)} & {(X \otimes U) \otimes Y}
            \arrow["\lambda", from=1-1, to=1-2]
            \arrow["{\lambda \otimes 1}"', from=1-1, to=2-2]
            \arrow["\alpha", from=1-2, to=2-2]
            \arrow["{\lambda \otimes 1}", from=1-4, to=1-5]
            \arrow["{1 \otimes \lambda}"', from=1-4, to=2-4]
            \arrow["{\beta \otimes 1}", from=1-5, to=2-5]
            \arrow["\alpha"', from=2-4, to=2-5]
        \end{tikzcd}\]
        \item If $(U', u')$ is another identity object in $(\~C, \otimes)$, there is a unique isomorphism $\vp: U \to U'$ making the following square commute:
        \[\begin{tikzcd}
            U & {U \otimes U} \\
            {U'} & {U' \otimes U'}
            \arrow["u", from=1-1, to=1-2]
            \arrow["\vp"', from=1-1, to=2-1]
            \arrow["{\vp \otimes \vp}", from=1-2, to=2-2]
            \arrow["{u'}"', from=2-1, to=2-2]
        \end{tikzcd}\]
    \end{enumerate}
\end{prop}

There is another natural isomorphism called the \emph{right unitor}, defined by
\[\rho_X := \alpha_{UX} \circ \lambda_X: X \to X \otimes U\]
it satisfies analagous properties to $\lambda$, but with $U$ on the right-hand side of the tensor product. The existence of $\lambda$ and $\rho$ tells us that every identity object in $(\~C, \otimes)$ is determined up to a unique isomorphism, so we may speak of \emph{the} identity object in $(\~C, \otimes)$. We typically denote this object by $(\mathbf{1}, e)$.

\begin{rmk}
    We can extend the functor $\otimes: \~C \times \~C \to \~C$ to a functor
\[\bigotimes_{i = 1}^n: \~C^n \to \~C\]
which is determined up to a unique isomorphism.
\end{rmk}

\begin{defn}
An object $L$ in a tensor category $(\~C, \otimes)$ is called \emph{invertible} if the functor
\begin{align*}
    L \otimes (-): \~C &\to \~C \\ 
    X & \mapsto L \otimes X
\end{align*}
is an equivalence of categories. Thus $L$ is invertible if and only if there exists an object $L^{-1} \in \~C$ such that $L \otimes L^{-1} = \mathbf{1}$. An \emph{inverse} of $L$ is a pair $(L^{-1}, \delta)$ where $\delta$ is an isomorphism
\[\delta: L \otimes L^{-1} \isom \mathbf{1}.\]
This definition is symmetric: $(L, \delta)$ is also an inverse of $L^{-1}$.
\end{defn}
Invertible objects are determined up to a unique isomorphism: If $(L_1, \delta_1)$ and $(L_2, \delta_2)$ are both inverses of $L$, there is a unique isomorphism $\psi: L_1 \to L_2$ such that $\delta_1 = \delta_2 \circ (1 \otimes \psi)$ 
\[\delta_1 = L \otimes L_1 \To{1 \otimes \psi} L \otimes L_2 \To{\delta_2} \mathbf{1}.\]

\subsubsection{Internal Hom and rigidity}
Let $(\~C, \otimes)$ be a tensor category. Suppose the functor 
    \[
        \~C^\op \to \~C, \quad T \mapsto \~C(T \otimes X, Y)
    \]
is representable, and denote by $\underline{\Hom}(X, Y)$ the representing object for this functor. Recall this means that for all $T \in \~C$, we have a natural isomorphism
\[\~C(T \otimes X, Y) \cong \~C(T, \underline{\Hom}(X,Y)).\]
The isomorphism is governed by the map
\[\ev_{XY}: \underline{\Hom}(X, Y) \otimes X \to Y\]
corresponding to $\id_{\underline{\Hom}(X, Y)}$. Each morphism $g: T \otimes X \to Y$ corresponds to a unique morphism $f: T \to \underline{\Hom}(X, Y)$ such that $e_{XY} \circ (f \otimes \id_X) = g$. This situation is depicted below:
\[\begin{tikzcd}
    T && {T\otimes X} \\
    {\underline{\Hom}(X, Y)} && {\underline{\Hom}(X, Y) \otimes X} && Y
    \arrow["f"', dashed, from=1-1, to=2-1]
    \arrow["{f \otimes \id}"', dashed, from=1-3, to=2-3]
    \arrow["g", from=1-3, to=2-5]
    \arrow["\ev"', from=2-3, to=2-5]
\end{tikzcd}\]

\begin{defn}
    $\underline{\Hom}(X, Y)$ is called an \emph{internal Hom object}. From this point forward, we assume $\underline{\Hom}(X, Y)$ exists for all $X, Y \in \~C$.
\end{defn}

\begin{ex}
    In $\Mod_R$, $\underline{\Hom}(X, Y)$ is just given by the $R$-module $\Hom_R(X, Y)$. This is a consequence of the familiar isomorphism
    \[\Hom_R(T, \Hom_R(X, Y)) \cong \Hom_R(T \otimes_R X, Y)\]
    The map $\ev_{XY}$ is given by the map
    \begin{align*}
        \ev_{XY}: \Hom_R(X, Y) \otimes X &\to Y \\ 
        f \otimes x &\mapsto f(x)
    \end{align*} 
    (this explains the notation $\ev_{XY}$).
\end{ex}

Given $X, Y, Z \in \~C$, we have a composition map
\[\underline{\Hom}(X, Y) \otimes \underline{\Hom}(Y, Z) \to \underline{\Hom}(X,Z),\]
which is given explicitly by the composite
\[\iHom(X, Y) \otimes \iHom(Y, Z) \To{\ev} \iHom(Y, Z) \otimes Y \To{\ev} Z.\]
Moreover, we have an isomorphism
\begin{equation}
    \iHom(Z, \iHom(X,Y)) \isom \iHom(Z \otimes X, Y),
\end{equation}
which further induces an isomorphism
\begin{align*}
    \~C(W, \iHom(Z, \iHom(X, Y))) &\cong \~C(W \otimes Z, \iHom(X, Y)) \\ 
                                  &\cong \~C(W \otimes Z \otimes X, Y) \\ 
                                  &\cong \~C(W, \iHom(Z \otimes X, Y)).
\end{align*}
An immediate consequence of this is the following:
\begin{equation}
    \~C(\mathbf{1}, \iHom(X, Y)) \cong \~C(\textbf{1} \otimes X, Y) = \~C(X, Y).
\end{equation}

\begin{defn}
    The \emph{dual} of an object $X\in \~C$ is $V^\vee := \iHom(X, \mathbf{1})$. 
\end{defn}

We therefore obtain a map
\[\ev_X : X^\vee \otimes X \to \mathbf{1}\]
and an isomorphism
\begin{equation}
    \~C(W, \iHom(X, \mathbf{1})) = \~C(W, X^\vee) \cong \~C(W \otimes X, \mathbf{1}) \quad \text{for all } W, X \in \~C.
\end{equation}
Moreover, the assignment $X \mapsto X^\vee$ is functorial: given a morphism $f: X \to Y$, we define a morphism $^{t}f: X^\vee \to Y^\vee$ as the unique arrow $Y^\vee \to X^\vee$ making the following square commute:
\[\begin{tikzcd}
    {Y^\vee \otimes X} & {X^\vee \otimes X} \\
    {Y^\vee \otimes Y} & {\mathbf{1}}
    \arrow["{^{t}f \otimes \id}", from=1-1, to=1-2]
    \arrow["{\id \otimes f}"', from=1-1, to=2-1]
    \arrow["{\ev_X}", from=1-2, to=2-2]
    \arrow["{\ev_Y}"', from=2-1, to=2-2]
\end{tikzcd}\]
this gives a functor
\[(-)^\vee: \~C^\op \to \~C.\]
When $f: X \to Y$ is an isomorphism, we define 
\[f^\vee:=  (^{t}f)^{-1}: X^\vee \to Y^\vee,\]
so that
\[\ev_Y \circ (f^\vee \otimes f) = \ev_X: X^\vee \otimes X \to \mathbf{1}.\]
\begin{defn}
    Let $X$ be an object in a tensor category $\~C$, and let $i_X: X \to X^{\vee \vee}$ be the morphism corresponding to $\ev_X \circ \psi: X \otimes X^\vee \to \mathbf{1}$. $X$ is said to be \emph{reflexive} if $i_X$ is an isomorphism.
\end{defn}

\begin{defn}
    We say that a tensor category $(\~C, \otimes)$ is \emph{rigid} if 
    \begin{enumerate}[(i)]  
        \item $\iHom(X, Y)$ exists for all $X, Y \in \~C$. 
        \item The canonical morphism
        \[\iHom(X_1, Y_1) \otimes \iHom(X_2, Y_2) \to \iHom(X_1 \otimes X_2, Y_1 \otimes Y_2)\]
        is an isomorphism for all $X_1, X_2, Y_1, Y_2 \in \~C$.
        \item Every object in $\~C$ is reflexive.
    \end{enumerate}
\end{defn}
\subsubsection{Tensor functors} We now define a notion of functors between tensor categories:
\begin{defn}
    If $(\~C, \otimes)$ and $(\~C', \otimes')$ are tensor categories, a \emph{tensor functor} $(\~C, \otimes) \to (\~C', \otimes')$ is a pair $(F, c)$ where $F: \~C \to \~C'$ is a functor and $c_{XY}: F(X) \otimes F(Y) \to F(X \otimes Y)$, such that the following properties are satisfied:
    \begin{enumerate}
        \item For all $X, Y, Z \in \~C$, we have a commuting diagram
        \[\begin{tikzcd}
    {FX \otimes (FY \otimes FZ)} & {FX \otimes F(Y \otimes Z)} & {F(X \otimes(Y \otimes Z))} \\
    {(FX \otimes FY) \otimes FZ} & {F(X\otimes Y) \otimes FZ} & {F((X \otimes Y) \otimes Z)}
    \arrow["{\id \otimes c}", from=1-1, to=1-2]
    \arrow["{\alpha'}"', from=1-1, to=2-1]
    \arrow["c", from=1-2, to=1-3]
    \arrow["{F(\alpha)}", from=1-3, to=2-3]
    \arrow["{c \otimes \id}"', from=2-1, to=2-2]
    \arrow["c"', from=2-2, to=2-3]
\end{tikzcd}\]
\item The following square commutes for all $X, Y \in \~C$:
\[\begin{tikzcd}
    {FX \otimes FY} & {F(X\otimes Y)} \\
    {FY \otimes FX} & {F(Y \otimes X)}
    \arrow["c", from=1-1, to=1-2]
    \arrow["{\beta'}"', from=1-1, to=2-1]
    \arrow["{F(\beta)}", from=1-2, to=2-2]
    \arrow["c"', from=2-1, to=2-2]
\end{tikzcd}\]
\item If $(U, u)$ is an identity object in $\~C$, then $(F(U), F(u))$ is an identity object in $\~C'$.
    \end{enumerate}
When $F: \~C \to \~C'$ is an equivalence of categories, we say that $(F, c)$ is an \emph{equivalence of tensor categories} or \emph{tensor equivalence}.
\end{defn}

\begin{defn}
    Suppose $(F, c)$ and $(G, d)$ are tensor functors $\~C \to \~C'$. A \emph{morphism of tensor functors} $(F, c) \to (G, d)$ is a natural transformation $\lambda: F \to G$ such that the following diagram commutes for all $X, Y \in \~C$:
\[\begin{tikzcd}
    {F(X)\otimes F(Y)} & {F(X \otimes Y)} \\
    {G(X) \otimes G(Y)} & {G(X \otimes Y)}
    \arrow["c", from=1-1, to=1-2]
    \arrow["{\lambda_X\otimes\lambda_Y}"', from=1-1, to=2-1]
    \arrow["{\lambda_{X \otimes Y}}", from=1-2, to=2-2]
    \arrow["d"', from=2-1, to=2-2]
\end{tikzcd}\]
This definition extends inductively to any finite family $(X_i)_{i=1}^n$ of objects in $\~C$, where we now require that the square
    \[\begin{tikzcd}
    {\bigotimes_{i=1}^n F(X_i)} & {F\left(\bigotimes_{i=1}^n X_i \right)} \\
    {\bigotimes_{i=1}^n G(X_i)} & {G\left(\bigotimes_{i=1}^n X_i \right)}
    \arrow["c", from=1-1, to=1-2]
    \arrow["{\otimes_{i=1}^n \lambda_{X_i}}"', from=1-1, to=2-1]
    \arrow["{\lambda_{\otimes_{i=1}^n X_i}}", from=1-2, to=2-2]
    \arrow["d"', from=2-1, to=2-2]
\end{tikzcd}\]
commutes. We write $\Fun^\otimes(\~C, \~C')$ for the category of tensor functors $(\~C, \otimes) \to (\~C', \otimes')$.
\end{defn}
When the tensor categories $\~C$ and $\~C'$ are both rigid, morphisms in $\Fun^\otimes(\~C, \~C')$ have a particularly nice description:
\begin{prop}[\cite{deligne1982tannakian}, Proposition 1.13]
    Let $(F, c)$ and $(G, d)$ be tensor functors $\~C \to \~C'$. If $\~C$ and $\~C'$ are rigid, then every morphism of tensor functors $\lambda: F \to G$ is an isomorphism.
\end{prop}

Recall that a \emph{subcategory} $\~B \mono \~C$ consists of a subclass $\ob(\~B)$ of $\ob(\~C)$, together with a subclass $\~B(X, Y)$ of $\~C(X,Y)$ (for all $X,Y \in \~B$) which is closed under composition and identities. A subcategory $\~B \mono \~C$ is said to be \emph{full} if for all $X, Y \in \~B $ we have $\~B(X,Y) = \~C(X, Y)$. A subcategory $\~B \mono \~C$ is \emph{strictly full} if it is full and closed under isomorphisms, that is, for every isomorphism $h \in \~C(X,Y)$ with $X \in \~B$, we have $h \in \~B(X,Y)$ (so $Y$ and $h^{-1}$ must also belong to $\~B$).
\begin{defn}
Let $\~B$ be a strictly full subcategory of a tensor category $(\~C, \otimes)$. We say that $\~B$ is a \emph{tensor subcategory} of $\~C$ if it is closed under finite tensor products. A subcategory $\~B$ of a rigid tensor category is called a \emph{rigid tensor subcategory} if $X^\vee \in \~B$ whenever $X \in \~B$.
\end{defn}

\subsubsection{Abelian tensor categories and examples} Recall that an \emph{additive category} is a category $\~A$ which is enriched over the category of abelian groups, so that $\~A(X, Y)$ is an abelian group for all $X, Y \in \~A$ and the morphisms between hom-sets are morphisms of abelian groups.

\begin{defn}
    An \emph{additive} (resp. \emph{abelian}) \emph{tensor category} is a tensor category $(\~C, \otimes)$ such that $\~C$ is additive (resp. abelian) and the functor $\otimes: \~C \times \~C \to \~C$ is bi-additive.
\end{defn} 

Suppose $(\~C, \otimes)$ is an additive tensor category and $(e, \mathbf{1})$ is an identity object. Then the endomorphisms of $\mathbf{1}$ form a ring $\End(\mathbf{1})$ which acts on each $X \in \~C$ via the left unitor map $\lambda_X: X \isom \mathbf{1} \otimes X$ of Proposition \ref{prop:LeftUnitor}. Moreover, $\End(\mathbf{1})$ is commutative since this action commutes with the endomorphisms of $X$. If we have another identity object $(\mathbf{1}', e')$, then the unique isomorphism $(\mathbf{1}, e) \isom (\mathbf{1}',e')$ of Proposition \ref{prop:LeftUnitor} gives an isomorphism $\End(\mathbf{1}) \cong \End(\mathbf{1}')$. This implies that each set of morphisms $\~C(X, Y)$ has the structure of $\End(\mathbf{1})$-module such that the composition operation $\circ$ is $\End(\mathbf{1})$-bilinear and $\otimes$ is bilinear.

\begin{prop}[\cite{deligne1982tannakian}, Proposition 1.16]
    Let $(\~C, \otimes)$ be a rigid tensor category. If $\~C$ is abelian, then $\otimes$ is bi-additive and commutes with direct and inverse limits in each variable. In particular, the functor $\otimes$ is exact in each variable.
\end{prop}
    \begin{ex}\hfill
        \begin{enumerate}
            \item As we have just seen, $\Rep_k(G)$ is a rigid abelian tensor category with $\End(\mathbf{1}) = k$. 

            \item The category $\Vec^\fd_k$ of finite-dimensional vector spaces over $k$ is a rigid abelian tensor category. Moreover, since any linear endomorphism $k \to k$ is determined by a scalar $\lambda \in k$, we have $\End(k) \cong k$.
            \item The category $\Mod^\fg_R$ of finitely-generated modules over a commutative ring $R$ is an abelian tensor category with $\End(\mathbf{1}) = \End(R) = R$, but it will not generally be rigid, since we may not have an isomorphism $M \cong M^{**}$ for every $R$-module $M$ (for instance, the $\Z$-module $\Z/2$ has $(\Z/2)^* = \Hom_\Z(\Z/2, \Z) \cong 0$, so $(\Z/2)^{**} \cong 0$).
            \item The category $\Proj_R^\fg$ of finitely generated projective modules over a commutative ring $R$ is a rigid additive tensor category with $\End(\mathbf{1}) = R$, but it is not necessarily abelian, since kernels may not always exist in this category. 
        \end{enumerate}
    \end{ex}
\begin{defn}
    A \emph{neutral Tannakian category} over $k$ is a pair $(\~C, \omega)$ where $\~C$ is a rigid abelian tensor category with $\End(\mathbf{1}) = k$ and $\omega: \~C \to \Vec^\fd_k$ is an exact faithful $k$-linear tensor functor, called the \emph{fibre functor}.
\end{defn}
\subsection{Tannaka duality} We are now ready to put all of the formalism of the previous section to good use. Suppose $G = \Spec(A)$ is an affine group scheme over $k$. We begin by reviewing 
% We now record some useful properties of coalgebras linear representations.
% \begin{prop}[\cite{deligne1982tannakian}, 2.3.-2.6.]
% \begin{enumerate}
%     \item Let $C$ be a $k$-coalgebra and $(V, \rho)$ a comodule over $C$. Then every finite subset of $V$ is contained in a sub-comodule of $V$ with finite dimension over $k$.
%     \item 
%     \item 
%     \item  
% \end{enumerate}
% \end{prop}
% \begin{proof}
%     To prove (1), let us suppose 

%     For (2), notice that the collection $\{W_i\}_{i \in I}$ of sub-comodules of $V$ forms a partially ordered set under inclusion. This set is also directed, since given two sub-comodules $W_1$ and $W_2$ of $V$, we can form their sum $W_1 + W_2$, which is by definition the smallest sub-comodule of $V$ containing $W_1$ and $W_2$. Moreover, the union of this directed system is clearly equal to $V$. Applying the correspondence in Proposition \ref{prop:DM2.2} gives the assertion in (2).
% \end{proof}

Let $(\~C, \omega)$ denote

\section{The Nori fundamental group scheme} 
\subsection{Preliminaries}
In this section we examine Nori's construction of the fundamental group scheme \cite{nori76}. Suppose $G$ is an affine group scheme over the field $k$. Let $R = k[G]$, and as in the previous sections, write $\Rep(G)$ for the category of finite-dimensional left representations of $G$, with trivial $G$-representation $\mathbf{1}_G$. Denote by $\omega: \Rep(G) \to \Vec_k^\fd$ the forgetful functor. Then $(\Rep(G), \otimes, \omega, \mathbf{1}_G)$ is a neutral Tannakian category; in particular, the following axioms are satisfied:
\begin{enumerate}
    \item[(T1)] $\Rep(G)$ is an abelian category over $k$.
    \item[(T2)] $\ob\left(\Rep(G)\right)$ is a set.
    \item[(T3)] $\otimes: \Rep(G) \times \Rep(G) \to \Rep(G)$ is a $k$-additive faithful exact functor.
    \item[(T4)] $\otimes: \Rep(G) \times \Rep(G) \to \Rep(G)$ is $k$-linear in each variable:
    \[\omega \circ \otimes = \otimes \circ (\omega \times \omega) \]
    \item[(T5)]\label{T5} $\otimes$ is associative and preserves $\omega$ in the following sense: Let $\alpha$ be associativity map of \autoref{ssub:2.1.1}. Given $X, Y, Z \in \~C$, $\omega(\alpha(X, Y, Z))$ gives an isomorphism $\omega X \otimes (\omega Y \otimes \omega Z) \cong (\omega X \otimes \omega Y) \otimes \omega Z$, and we require that this isomorphism coincides with the usual associativity map (i.e., the one which gives associativity of $\otimes$ in $\Vec_k$).
    \item[(T6)]\label{T6} $\otimes$ is commutative and preserves $\omega$ in the sense of (\hyperref[T5]{T5}) above.
    \item[(T7)]\label{T7} There is an object $\mathbf{1} \in \Rep(G)$ and an isomorphism $k \to \omega(\mathbf{1})$ preserving $\omega$ such that $\mathbf{1}$ is an identity object in $\Rep(G)$.
    \item[(T8)] For every object $X \in \Rep(G)$ such that $\omega(X)$ has dimension $1$, there is an object $X^\vee \in \Rep(G)$ and an isomorphism $X \otimes X^\vee \isom \mathbf{1}$.
\end{enumerate}
Nori \cite{nori76} defines a neutral Tannakian category\footnote{Nori actually uses the term ``Tannaka category'', but we will stick to the language of neutral Tannakian categories.} as any tuple $(\~C, \otimes, \omega, \mathbf{1})$ satisfying these conditions. We will adopt this more object-focused approach in the coming sections. 
\begin{defn} Let $\~C$ be a category satisfying (\hyperref[T2]{T2}) and (\hyperref[T3]{T3}), and let $S$ be a set of objects in $\~C$. We write
    \[\overline{S} := \left\{X \in \~C: \exists P_i \in S, 1 \le i \le t, \text{ and } V_1, V_2 \in \~C \text{ such that } V_1 \subseteq V_2 \subseteq \bigoplus_{i=1}^t P_i \text{ and } W \cong V_2/V_1 \right\}.\]
Denote by $\~C(S)$ the full subcategory of $\~C$ with $\ob \left(\~C(S)\right) = \overline{S}$. It is an abelian category. $S$ is said to \emph{generate} $\~C$ if $\ob(\~C) = \overline{S}$.
\end{defn}

For a ringed space $(X, \~O_X)$, recall that a sheaf $\~F$ of $\~O_X$-modules is said to be of \emph{finite type} if for each $x \in X$ there is an open neighbourhood $U \subseteq X$ of $x$ such that $\~F|_U$ is generated by finitely many sections. That is, for every $x \in X$ there is an open neighbourhood $U$ of $x$ and a surjection of $\~O$-modules
    \[\~O_U^{\oplus n} \epi \~F|_U\] 
(where $\~O_U = \~O_X|_U$). Hence if $\~F$ has finite type, then for any $x \in X$ there is some neighbourhood $U \subseteq X$ and a collection of sections $s_1, \hdots, s_n \in \~F(U)$ such that for all $V \subseteq U$ and $t \in \~F(V)$, we can write
\[t =\sum_{i = 1}^n a_i \left(s_i|_V\right), \quad a_i \in \~O_X(V).\]
Recall that a \emph{coherent sheaf} is a sheaf $\~F$ of $\~O_X$-modules which has finite type, such that for any collection of sections $s_1, \hdots, s_n \in \~F(U)$, the kernel of the associated map
    \[\~O_U^{\oplus n} \to \~F|_U, \quad (a_1, \hdots, a_n) \mapsto \sum_{i=1}^n a_i s_i\]
has finite type. For each coherent sheaf $\~F$ on $X$, there is a surjection 
\[\~O_U^{\oplus n} \epi \~F|_U.\]
If $\~K$ denotes the kernel of this map, then we have an exact sequence
\[0 \to \~K \mono \~O_U^{\oplus n} \epi \~F|_U \to 0,\]
but also $\~K$ has finite type, so there is another surjection $\~O_U^{\oplus m} \epi \~K$. This gives a right-exact sequence of sheaves:
\begin{equation}\label{eq:1.1}
    \~O_U^{\oplus m} \to \~O_U^{\oplus n} \epi \~F|_U \to 0.
\end{equation}
Hence $\~F|_U$ is the cokernel of the map $\~O_U^{\oplus m} \to \~O_U^{\oplus n}$. Hence every coherent sheaf is the cokernel of a morphism of sheaves $\~O_X^{\oplus m} \to \~O_X^{\oplus n}$. The exactness of (\ref{eq:1.1}) is often used as the \emph{definition} of a coherent sheaf. We say that a sheaf $\~F$ of $\~O_X$ modules on a ringed space $(X, \~O_X)$ is \emph{quasi-coherent} if for every point $x \in X$ there is a neighbourhood $U$ of $x$ and an exact sequence
\begin{equation}\~O^{\oplus I}_U \to \~O^{\oplus J}_U \epi \~F|_U \to 0\end{equation}
The category of quasi-coherent sheaves on $X$ is denoted $\QCoh_X$. We will mostly work with quasi-coherent sheaves in these notes.

Throughout these notes, we will also make frequent reference to \emph{flat} $R$-modules. A module  $M$ over a ring $R$ is called \emph{flat} if the tensor product functor $(-)\otimes_R M$ is exact, and a ring morphism $R \to S$ is \emph{flat} if $S$ is a flat $R$-module under the action induced by this homomorphism. A morphism $f: X \to Y$ of ringed spaces is \emph{flat at $x \in X$} if the associated ring homomorphism $\~O_{Y, f(x)} \to \~O_{X, x}$ is flat, and we say that $f$ is \emph{flat} if it is flat at every $x \in X$. When $\~F$ is an $\~O_X$-module, we say that $\~F$ is \emph{flat} if the functor 
\[
   (-)\otimes_{\~O_X}\~F: \Mod_{\~O_X} \to \Mod_{\~O_X}, \quad \~G \mapsto \~G \otimes_{\~O_X} \~F
\]
is exact (for more details, see \S3.9 in \cite{hartshorne}). 
\subsection{Principal bundles}
Let $X$ be a scheme over $k$ and $G$ an affine group scheme over $k$. In this section we introduce \emph{principal $G$-bundles}. We will make use of the following definition from Nori's paper \cite{nori76}:

\begin{defn}
    A morphism $j: P \to X$ is called a \emph{principal $G$-bundle} on $X$ if the following conditions are satisfied:
    \begin{enumerate}
        \item $j$ is a surjective flat affine morphism. 
        \item The map $\Phi: P \times G \to P$ gives an action of $G$ on $P$ such that $j \circ \Phi = j \circ p_1$, i.e., the square
        \[\begin{tikzcd}
            {P\times G} & P \\
            P & X
            \arrow["{p_1}", from=1-1, to=1-2]
            \arrow["\Phi"', from=1-1, to=2-1]
            \arrow["j", from=1-2, to=2-2]
            \arrow["j"', from=2-1, to=2-2]
        \end{tikzcd}\]
        commutes (where $p_1$ is the projection onto the first factor).
        \item The map $\Psi := (p_1, \Phi): P \times G \to P \times_X P$ depicted in the following pullback diagram is an isomorphism:
        \[\begin{tikzcd}[column sep = small]
    {P \times G} \\
    & {P\times_X P} & P \\
    & P & X
    \arrow["\Psi"{description}, dashed, from=1-1, to=2-2]
    \arrow["{p_1}", curve={height=-6pt}, from=1-1, to=2-3]
    \arrow["\Phi"', curve={height=6pt}, from=1-1, to=3-2]
    \arrow[from=2-2, to=2-3]
    \arrow[from=2-2, to=3-2]
    \arrow["\lrcorner"{anchor=center, pos=0.125}, draw=none, from=2-2, to=3-3]
    \arrow["j", from=2-3, to=3-3]
    \arrow["j"', from=3-2, to=3-3]
\end{tikzcd}\]
    \end{enumerate}
\end{defn}
Now suppose $V$ is a finite-dimensional representation of $G$ and we are given a principal $G$-bundle $j: P \to X$. We therefore have a right action of $G$ on $P$ and a left action of $G$ on $V$. $G$ now acts on the trivial vector bundle $P \times V$ by $g \cdot (p, v) = (pg, g^{-1}v)$. In this situation we define a new vector bundle $j^G: P \times^G V \epi X$ via the quotient under this action:
\[P \times^G V := (P \times V)/G.\]
We can now define a sheaf $F_PV$ on $X$ by taking sections of this bundle: For any open set $U \subseteq X$, define
\[F_PV(U) := \Gamma\left(P \times^G V, U\right) = \{s: U \to P \times^G V : s \circ j^G = \id_U\}.\]
Recall that the sheaf of sections of a vector bundle is quasicoherent since any vector bundle has a local trivialization. In particular, the sheaf $F_PV$ defined above is quasicoherent, so we have a functor
\begin{align*}
    F_P: \Rep(G) &\to \QCoh_X\\
    V &\mapsto F_PV
\end{align*}
This functor has the following properties:
\begin{enumerate}
    \item[(F1)]\label{F1} $F_P$ is exact and $k$-additive.
    \item[(F2)]\label{F2} $F_P \circ \otimes = \otimes \circ (F_P \times F_P)$.
    \item[(F3)] \label{F3}The analogous statements to (\hyperref[T5]{T5}), (\hyperref[T6]{T6}), and (\hyperref[T7]{T7}) hold.
    \item[(F4)] \label{F4}If $\rank(V) = n$, then $F_PV$ is a locally free sheaf of rank $n$, so that $F_P$ is faithful.
\end{enumerate}
Unless otherwise specified, $F$ denotes a functor satisfying (\hyperref[F1]{F1})-(\hyperref[F4]{F4}).

\begin{lem}
    Let $\Rep'(G)$ denote the category of all (possible infinite-dimensional) representations of $G$. Given a functor $F$ as above, there is a unique functor $\overline{F}: \Rep'(G) \to \QCoh(X)$, defined for $V \in \Rep'(G)$ by
    \[\overline{F}(V) := \colim_{W \subseteq V}FW\]
    where the colimit is over all finite-dimensional $G$-invariant subspaces $W \subseteq V$. Moreover, this functor has the following properties:
    \begin{enumerate}
        \item $\overline{F}$ satisfies (\hyperref[F1]{F1}), (\hyperref[F2]{F2}), and (\hyperref[F3]{F3}).
        \item $\overline{F}|_{\Rep(G)} = F$.
        \item $\overline{F}V$ is flat for all $V \in \Rep(G)'$, and faithfully flat when $V \ne 0$.
        \item $\overline{F}$ preserves direct limits.
    \end{enumerate}
\end{lem}

\begin{proof}
    (1) is just a simple check of the axioms (\hyperref[F1]{F1})-(\hyperref[F3]{F3}). For example, (\hyperref[F1]{F1}) follows from the fact that filtered colimits preserve exactness in an abelian category, and (\hyperref[F2]{F2}) boils down to showing that the tensor product commutes with filtered colimits. For (2), notice that if $V$ is finite-dimensional, then we are taking the colimit over a filtered category with a terminal object (namely $V$), so $\overline{F}(V) = \colim_{W \subseteq V}FW = F(V)$. For (3), recall that any free module is flat, hence any locally free sheaf is flat, so each sheaf $FW$ in the colimit defining $\overline{F}V$ is flat. Since a filtered colimit of flat modules (or sheaves) is flat (see \cite{stacks}, Lemma 03EU), we deduce (3). (4) is obvious. 
\end{proof}

\begin{notation}
    From this point forward, we set $\overline{F} = F$.
\end{notation}
\newpage
\printbibliography
\end{document}


